% Generated by scripts/assemble.py; corrections belong in source/corrections.json.
\examhead{37}{华东师范大学 2025 年数学分析试卷(回忆版)}{40}
% p041-b03; source PDF page 41
\par\Needspace{5\baselineskip}
\blockanchor{p041-b03}
\majorquestion{一}填空题.\par
% p041-b04; source PDF page 41
\blockanchor{p041-b04}
\question{1}$\lim_{n\to\infty}\sqrt[n]{\left(1+\frac{1}{2}\right)\left(1+\frac{2}{3}\right)\cdots\left(1+\frac{n-1}{n}\right)}=$  \_\_\_\_.\par
% p041-b05; source PDF page 41
\blockanchor{p041-b05}
\question{2}$\int_0^\pi \frac{x\sin^2x}{1 + \cos^2x}\mathrm{d}x =$\par
% p041-b06; source PDF page 41
\blockanchor{p041-b06}
\question{3}$\lim_{x\to\infty}\left(\cos\frac{a}{x}\right)^{x^{2}}=$  \_\_\_\_.\par
% p041-b07; source PDF page 41
\blockanchor{p041-b07}
\question{4}$\sum_{n=1}^{\infty} (-1)^{n-1} \frac{\sin nx}{n}$ 是 \_\_\_\_ 在 $(- \pi, \pi)$ 上的傅里叶级数.\par
% p041-b08; source PDF page 41
\blockanchor{p041-b08}
\question{5}求和函数 $\sum_{n=0}^{\infty} \frac{n^{2}+1}{n!2^{n}} x^{n} = \underline{\quad}$.\par
% p041-b09; source PDF page 41
\blockanchor{p041-b09}
\question{6}由柱面坐标表示的曲面 $z = z(r, \theta)$, 其上任意一条光滑曲线的弧长微分 $\mathrm{ds}$ 满足 $\mathrm{ds}^2 = \_\_\_\_ \mathrm{dr}^2 + \_\_\_\_ \mathrm{dr} \mathrm{d}\theta + \_\_\_\_ \mathrm{d}\theta^2$.\par
% p041-b10; source PDF page 41
\blockanchor{p041-b10}
\question{7}求第二型曲面积分 $\iint_{S} \frac{a^2 x \, \mathrm{d}y \, \mathrm{d}z + (z + a)^2 \, \mathrm{d}x \, \mathrm{d}y}{(x^2 + y^2 + z^2)^{\frac{3}{2}}} =$，其中 $S$ 为 $z = -\sqrt{a^2 - x^2 - y^2}$，取上侧.\par
% p041-b11; source PDF page 41
\par\Needspace{5\baselineskip}
\blockanchor{p041-b11}
\majorquestion{二}简答题.\par
% p041-b12; source PDF page 41
\blockanchor{p041-b12}
\question{1}解答如下问题:\par
% p041-b13; source PDF page 41
\blockanchor{p041-b13}
(1) $f(x)$ 为定义在 $(a, b)$ 上的函数, 若对任意的 $x_0 \in (a, b)$, 有 $\lim_{x \to x_0} f(x)$ 存在, 任取 $\{a_n\} \subset (a, b)$ 为 Cauchy 列, 那么是否有 $\{f(a_n)\}$ 为 Cauchy 列?\par
% p041-b14; source PDF page 41
\blockanchor{p041-b14}
(2) $f(x)$ 为定义在 $(a, b)$ 上的函数, 若对 $(a, b)$ 中任意的 Cauchy 列 $\{a_n\}$, $\{f(a_n)\}$ 也是 Cauchy 列, 证明: $f(x)$ 在 $(a, b)$ 上一致连续.\par
% p041-b15; source PDF page 41
\blockanchor{p041-b15}
\question{2}设 $f(x)$ 在 $[a, b]$ 上可微.\par
% p041-b16; source PDF page 41
\blockanchor{p041-b16}
(i)  $f'(x)$  在  $[a, b]$  上连续.\par
% p041-b17; source PDF page 41
\blockanchor{p041-b17}
(ii) 对任意的 $\varepsilon > 0$, 存在 $\delta > 0$, 使得当 $0 < |h| < \delta$ 时, 有 $\left|\frac{f(x + h) - f(x)}{h} - f'(x)\right| < \varepsilon$ 对任意的 $x \in [a, b]$ 成立.\par
% p041-b18; source PDF page 41
\blockanchor{p041-b18}
. 讨论 (i) 能否推 (ii)? (ii) 能否推 (i)?\par
% p041-b19; source PDF page 41
\blockanchor{p041-b19}
\question{3}设 $\{f_n(x)\}$ 为 $[a, b]$ 上连续的函数列, 且 $\sum_{n=1}^{\infty} f_n(x)$ 在 $(a, b)$ 上一致收敛.\par
% p041-b20; source PDF page 41
\blockanchor{p041-b20}
(1) 证明: $\sum_{n=1}^{\infty} f_n(a), \sum_{n=1}^{\infty} f_n(b)$ 收敛.\par
% p041-b21; source PDF page 41
\blockanchor{p041-b21}
(2) 证明: $\sum_{n=1}^{\infty} f_n(x)$ 在 $[a, b]$ 上一致收敛.\par
% p041-b22; source PDF page 41
\blockanchor{p041-b22}
\question{4}设 $f(x)$ 为 $(- \infty, + \infty)$ 上的函数, 且在任意有限区间上可积, 若 $\int_{-\infty}^{+\infty}|f(x)|\mathrm{d}x$ 收敛, 证明: $F(x) = \int_{-\infty}^{+\infty}f(t)\cos 2\pi tx\mathrm{d}t$ 在 $(- \infty, + \infty)$ 上一致连续.\par
% p042-b00; source PDF page 42
\blockanchor{p042-b00}
\question{5}解答如下问题:\par
% p042-b01; source PDF page 42
\blockanchor{p042-b01}
(1) 对于 $\sum_{n=1}^{\infty} a_n$, 若对任意子列 $\{a_{n_k}\}$, 均有 $\sum_{k=1}^{\infty} a_{n_k}$ 收敛, 那么 $\sum_{n=1}^{\infty} a_n$ 是否绝对收敛?\par
% p042-b02; source PDF page 42
\blockanchor{p042-b02}
(2) 对于任意正整数 $i, j$, 均有 $\sum_{n=1}^{\infty} a_{i+nj}$ 收敛, 是否有 $\sum_{n=1}^{\infty} a_n$ 绝对收敛?\par
% p042-b03; source PDF page 42
\par\Needspace{6\baselineskip}
\blockanchor{p042-b03}
\question{6}设 $V, U$ 分别是 $\mathbb{R}^m, \mathbb{R}^n$ 上的开集, 定义向量函数\par
\nopagebreak[4]
% p042-b04; source PDF page 42
\blockanchor{p042-b04}
\displayformula{\boldsymbol {f}: V \mapsto U, (x _ {1}, \dots , x _ {m}) \rightarrow (f _ {1} (x _ {1}, \dots , x _ {m}), \dots , f _ {n} (x _ {1}, \dots , x _ {m})).}
% p042-b05; source PDF page 42
\par\Needspace{6\baselineskip}
\blockanchor{p042-b05}
称 f 在点  $x_{0}$  可微, 如果存在  $m \times n$  矩阵 A, 满足\par
\nopagebreak[4]
% p042-b06; source PDF page 42
\blockanchor{p042-b06}
\displayformula{\boldsymbol {f} (\boldsymbol {x} _ {0} + \boldsymbol {h}) = f (\boldsymbol {x} _ {0}) + \boldsymbol {h} A + o (\| \boldsymbol {h} \|).}
% p042-b07; source PDF page 42
\blockanchor{p042-b07}
其中  $\boldsymbol{h}=(h_{1},\cdots,h_{m}),o(\|\boldsymbol{h}\|)$  为各元素均为  $\|h\|$  高阶无穷小的 m 维行向量. 此外, 定义 Jacobi 阵为  $D_{x_{0}}^{f}=\left(\frac{\partial f_{j}}{\partial x_{i}}\right)_{m\times n}\bigg|_{x=x_{0}}$ .\par
% p042-b08; source PDF page 42
\blockanchor{p042-b08}
(1) 证明: 若 f 在  $x_{0}$  处可微, 则  $D_{x_{0}}^{f}=A$ .\par
% p042-b09; source PDF page 42
\blockanchor{p042-b09}
(2) 进一步, 若 $g$ 在 $f(\pmb{x}_0)$ 处可微, 则 $g \circ f$ 在 $\pmb{x}_0$ 处可微, 且有 $A = D_{f(\pmb{x}_0)}^g D_{\pmb{x}_0}^f$.\par
% p042-b10; source PDF page 42
\blockanchor{p042-b10}
(3) 未知...\par
